\section{问题一：投放策略的合理性与时间规律}
\label{sec:q1}
本节先说明数据口径与质量，再从设计质量与创意、关键词管理、出价与预算、投放时间四个方面给出可检验的事实，最后估计注册归因模型与单元级弹性，为后续问题提供结构参数。

\subsection{数据说明与质量}
\label{sec:q1data}
所有数据经契约校验（列类型、取值范围、漏斗一致性：点击 $\le$ 展现、上方位展现 $\le$ 展现、首位展现 $\le$ 上方位展现、上方位点击 $\le$ 点击、上方位消费 $\le$ 消费）全部通过。Sheet1 有 \valUnitDays{} 条单元--日记录，覆盖 \valPlansCount{} 个方案、\valUnitsCount{} 个推广单元和 2025 年全部 365 天，（日期、方案、单元）无重复；各单元的投放起止不同（四个单元自 5 月 21 日起投，一个单元 11 月 28 日停投），春节等假期多数单元停投，\valZeroSpendUnitDays{} 条记录有展现而消费为零。全年合计消费 \valTotalSpend{} 元（题述“142 余万元”）、点击 \valTotalClicks{} 次、展现 \valTotalImpressions{} 次。Sheet2 的新注册数合计 \valTotalRegs{} 人。Sheet3 有 \valKeywordRows{} 条关键词记录（题述“6000 多个关键词”多于附件行数，本文以附件为准，假设~\ref{as:data}），关键词 ID 为匿名整数，共 \valUniqueKeywordIds{} 个不同 ID，其中 \valSharedKeywordIds{} 个 ID 被 2--4 个单元共用（涉及 \valSharedKeywordRows{} 条记录）；\valZeroKeywords{} 条（\valZeroKeywordSharePct\%）全年零消费零点击，其跳出率与时长以“/”标记；跳出率为 0--1 的小数，平均访问时长为“HH:MM:SS”或数值 0（84 行）。Sheet3 分单元的消费与点击合计与 Sheet1 完全对账（总消费差异 0.02 元，为四舍五入）。

\subsection{指标体系与四个方面的事实}
\label{sec:q1kpi}
以单元 $u$ 的年度合计定义漏斗指标
\begin{equation}\label{eq:kpi}
\mathrm{CTR}_u=\frac{\sum_t C_{u,t}}{\sum_t I_{u,t}},\qquad
\mathrm{CPC}_u=\frac{\sum_t S_{u,t}}{\sum_t C_{u,t}},\qquad
p^{\mathrm{top}}_u=\frac{\sum_t I^{\mathrm{top}}_{u,t}}{\sum_t I_{u,t}},\qquad
p^{1}_u=\frac{\sum_t I^{1}_{u,t}}{\sum_t I_{u,t}},
\end{equation}
关键词消费的集中度用按消费降序的 Lorenz 曲线与 Gini 系数度量：对有消费的 $n$ 个词按 $S_{(1)}\le\dots\le S_{(n)}$ 排序，
\begin{equation}\label{eq:gini}
G=\frac{n+1-2\sum_{k=1}^{n}\bigl(\sum_{j\le k}S_{(j)}\bigr)/\sum_{j}S_{(j)}}{n}.
\end{equation}
表~\ref{tab:unit-kpis} 给出各单元的年度 KPI，图~\ref{fig:unit-lorenz} 给出单元消费占比、CPC--CTR 散点与关键词 Lorenz 曲线。

\papertable{tab_unit_kpis}{各推广单元 2025 年投放 KPI（按消费额降序；CTR、CPC、上方位与首位占比按式~\eqref{eq:kpi}）}{tab:unit-kpis}

\begin{figure}[htbp]\centering
\begin{minipage}[t]{0.62\linewidth}\centering\includegraphics[width=\linewidth]{fig_unit_kpis.pdf}\end{minipage}\hfill
\begin{minipage}[t]{0.36\linewidth}\centering\includegraphics[width=\linewidth]{fig_lorenz.pdf}\end{minipage}
\caption{左：各推广单元年度消费额占比；中：单元 CPC--CTR 散点（对数纵轴，气泡面积 $\propto$ 消费占比，标签为单元 ID 后四位）；右：关键词消费额与点击量的 Lorenz 曲线（按消费额降序累计），虚线为均匀分布。}
\label{fig:unit-lorenz}
\end{figure}

\paragraph{设计质量与创意}
账户整体 CTR 为 \valOverallCtrPct\%，CPC \valOverallCpc{} 元。上方位（前三名）展现只占 \valTopSharePct\%（首位 \valFirstSharePct\%），却贡献了 \valTopClickSharePct\% 的点击与 \valTopSpendSharePct\% 的消费：上方位 CTR 为 \valTopRegionCtrPct\%，非上方位仅 \valOtherRegionCtrPct\%，相差 \valTopOverOtherCtrRatio{} 倍，而上方位 CPC（\valTopCpc{} 元）只比非上方位（\valOtherCpc{} 元）高约 17\%。可检验的事实是展现结构严重偏向低 CTR 区域：非上方位以 \valOtherRegionImpSharePct\% 的展现只换来 \valOtherRegionClickSharePct\% 的点击。需要强调的是，位置拍卖的标准模型把 CTR 分解为位置效应与广告质量的乘积\cite{edelman2007,varian2007}，附件既无关键词日数据也无真实排名，本文\emph{无法}把上方位与非上方位的 CTR 差异归因于创意本身；只能说明在既定的位置结构下，把展现从低 CTR 区域移向上方区域的空间很大。单元 \valTopCtrUnitId{} 只有 \valTopCtrUnitKeywords{} 个关键词，CTR 高达 \valTopCtrUnitCtrPct\%、上方位占比 \valTopCtrUnitTopSharePct\%、CPC 仅 \valTopCtrUnitCpc{} 元，属品牌/导航型流量，说明品牌词的创意与匹配是充分的；相反，单元 \valMinCtrUnitId{} 的 CTR 仅 \valMinCtrUnitCtrPct\%（表~\ref{tab:unit-kpis}），其创意或关键词匹配需要重新审视。

\paragraph{关键词管理与运用}
关键词呈极端长尾：\valZeroKeywordSharePct\% 的词全年零消费零点击；有消费的 \valActiveKeywords{} 个词中，前 10 个词占消费 \valTopTenSpendSharePct\%，前 100 个占 \valTopHundredSpendSharePct\%，Gini 系数 \valGiniSpendActive{}（图~\ref{fig:unit-lorenz} 右）。\valSharedKeywordIds{} 个关键词 ID 同时出现在多个单元中，同一账户的不同单元会在同一查询上相互竞价，抬高自身 CPC；关键词库中约四成的词从未获得流量，说明词库缺乏定期清理与匹配方式调整。

\paragraph{出价策略与预算}
消费在单元间同样集中：单元 \valBiggestUnitId{} 占全年消费 \valBiggestUnitSpendSharePct\%，各单元 CPC 从 \valMinUnitCpc{} 元到 \valMaxUnitCpc{} 元不等（表~\ref{tab:unit-kpis}）。第~\ref{sec:q1attr} 节的点击--消费弹性合并估计为 \valPooledGamma（标准误 \valPooledGammaSe），显著小于 1，即边际收益递减；12 个单元的自身回归估计 $\hat\gamma_u\in[\valMinUnitGammaHat,\valMaxUnitGammaHat]$，其中 \valUnitsGammaHatBelowOne{} 个显著小于 1、\valUnitsGammaHatAboveOne{} 个（品牌型单元 \valGammaAboveOneUnitId）显著大于 1（见第~\ref{sec:q1attr} 节）。消费最大的单元 \valBiggestUnitId{} 的采用弹性最低（表~\ref{tab:attr}）而 CPC 最高，说明预算在该单元上已进入边际成本快速上升的区域。把同一天的预算在单元之间按边际收益均衡重新分配（用与第~\ref{sec:q3} 节相同的注水机制、单元级弹性异质、上限取历史份额的 $\mu$ 倍），在问题三的两个窗口上分别可再提升预期注册 \valQthreeCrossUnitGainPctFeb\% 与 \valQthreeCrossUnitGainPctAug\%（表~\ref{tab:q3-sens}）。需要说明的是，问题三与问题四的主模型按假设~\ref{as:budget} 把每个单元的预算锁定为其 2025 年当日实际消费，因此\emph{不}执行这种跨单元再分配；它作为反事实口径单独报告，也是第~\ref{sec:limits} 节列出的局限之一。

\paragraph{投放策略与时间}
图~\ref{fig:daily} 显示日消费、点击与注册的全年走势：三条序列高度同步（注册与消费的相关系数约 0.8），春节后 2--3 月出现全年高峰，6--7 月为低谷；每周内呈明显的“周一、周日低”模式，法定节假日（阴影）期间三者同时深跌。这些规律与假日效应在下一小节用回归量化。

\paperfigure[0.7\linewidth]{fig_daily_series.pdf}{2025 年日消费额、日点击量与日新注册数（细线为日值，粗线为 7 日滑动均值；阴影为法定节假日，竖虚线为调休上班日）。}{fig:daily}

\subsection{日历回归与假日效应}
\label{sec:q1cal}
对账户日总量 $y_t\in\{R_t,\ S_t,\ C_t,\ I_t\}$ 建立对数线性回归
\begin{equation}\label{eq:calreg}
\log y_t=\alpha+\sum_{w\ne\text{周三}}\beta_w D^{w}_t+\sum_{m\ne 1}\mu_m M^{m}_t+\delta_H H_t+\delta_A A_t+\delta_P P_t+\delta_Q Q_t+e_t,
\end{equation}
其中 $D^{w}_t$、$M^{m}_t$ 为周几与月份哑变量（基准为周三与 1 月），$H_t,A_t,P_t,Q_t$ 分别为法定节假日、调休上班日、假前一日、假后一日的指示变量（假设~\ref{as:calendar}）。系数 $\delta$ 的效应以 $100(e^{\delta}-1)\%$ 报告，标准误采用 MacKinnon--White 的 HC3 异方差稳健形式\cite{mackinnon1985}。对 CPC、CTR 与上方位占比以水平值代替对数。为把假日的需求侧效应与公司主动减投分离，再估计控制当日消费额的模型
\begin{equation}\label{eq:calreg-ctrl}
\log R_t=\alpha'+\theta\log S_t+\sum_{w}\beta'_w D^{w}_t+\sum_{m}\mu'_m M^{m}_t+\delta'_H H_t+\delta'_A A_t+\delta'_P P_t+\delta'_Q Q_t+e'_t,
\end{equation}
其中 $\theta$ 为注册对消费的弹性，$\delta'_H$ 是同样投入下假日仍然存在的注册变化。作为非参数对照，用 Mann--Whitney $U$ 检验\cite{mann1947}比较法定节假日（\valHolidayDays{} 天）与普通工作日（\valWorkdayDays{} 天）的日注册数分布。

\papertable{tab_calendar_effects}{日历回归~\eqref{eq:calreg} 的效应估计（相对周三 / 普通日，\%；HC3 稳健 95\% 置信区间与 $p$ 值）}{tab:calendar}

\paperfigure[0.76\linewidth]{fig_calendar_effects.pdf}{周几、法定节假日、调休、假前与假后对日新注册数、消费额与点击量的效应（相对周三 / 普通日，\%），横线为 HC3 稳健 95\% 置信区间。}{fig:calendar}

表~\ref{tab:calendar} 与图~\ref{fig:calendar} 给出估计结果。法定节假日的注册数效应为 $\valHolidayRegEffectPct\%$（95\% CI $[\valHolidayRegCiLow,\ \valHolidayRegCiHigh]$，$p=\num{\valHolidayRegPSci}$），消费额 $\valHolidaySpendEffectPct\%$、点击量 $\valHolidayClicksEffectPct\%$、展现量 $\valHolidayImpEffectPct\%$ 同样显著；Mann--Whitney 检验给出假日与工作日注册中位数 \valHolidayRegMedian{} 与 \valWorkdayRegMedian{} 人（比值 \valHolidayRegMedianRatio），$p=\num{\valHolidayRegMannWhitneyP}$，两种方法结论一致。价格与质量指标没有假日效应：CPC 的假日系数为 \valHolidayCpcEffect{} 元（$p=\valHolidayCpcP$），CTR 为 \valHolidayCtrEffectPp{} 个百分点（$p=\valHolidayCtrP$）；上方位占比在假日下降 \valHolidayTopEffectPp{} 个百分点（$p=\valHolidayTopP$），与减投后排位后移一致。周几效应显著（Wald 检验 $p=\num{\valRegWeekdayWaldP}$）：周一注册 $\valMondayRegEffectPct\%$、周日 $\valSundayRegEffectPct\%$，周六 $\valSaturdayRegEffectPct\%$（$p=\valSaturdayRegP$），周二至周五与周三无显著差异；调休上班日与普通日无差异（$\valAdjustedWorkdayRegEffectPct\%$，$p=\valAdjustedWorkdayRegP$）；假前一日、假后一日的点估计为负但因样本少不显著。

控制消费额的模型~\eqref{eq:calreg-ctrl} 把上述总效应分解：注册对消费的弹性 $\theta=\valRegSpendElasticity$（SE \valRegSpendElasticitySe），$R^2$ 由 \valRegCalendarRsq{} 升至 \valRegGivenSpendRsq；同样投入下假日注册仍低 $\valHolidayRegGivenSpendEffectPct\%$（95\% CI $[\valHolidayRegGivenSpendCiLow,\ \valHolidayRegGivenSpendCiHigh]$，$p=\valHolidayRegGivenSpendP$），分解在对数尺度上精确闭合：$\delta_H=\delta'_H+\theta\,\delta^{S}_H$，需求侧解释了对数跌幅的 \valHolidayDemandLogSharePct\%，其余 \valHolidaySupplyLogSharePct\% 来自公司假日主动减投 $\valHolidaySpendEffectPct\%$（百分点不可相加，故不能用 $24.6/73.7$ 计算份额）。周一（$\valMondayRegGivenSpendEffectPct\%$，$p=\valMondayRegGivenSpendP$）与周日（$\valSundayRegGivenSpendEffectPct\%$，$p=\valSundayRegGivenSpendP$）在控制消费后不再显著，其点估计的需求侧跌幅远小于实际减投幅度。

假日减投是否合理，可用边际收益均衡原则判断。设注册对消费的响应为 $R=K S^{\theta}$，假日需求侧乘子为 $e^{\delta'_H}$，则假日与普通日边际注册相等要求 $e^{\delta'_H}\theta K S_H^{\theta-1}=\theta K S_N^{\theta-1}$，即最优的假日/普通日消费比
\begin{equation}\label{eq:holiday-ratio}
\frac{S_H^{*}}{S_N^{*}}=\exp\!\Bigl(\frac{\delta'_H}{1-\theta}\Bigr)=\valHolidayOptimalSpendRatioPct\%,
\end{equation}
即最优减投幅度的点估计为 \valHolidayOptimalCutPct\%。

式~\eqref{eq:holiday-ratio} 的两个输入 $\delta'_H$ 与 $\theta$ 来自同一个回归，必须联合传播其抽样不确定性：对二者的 HC3 协方差作 $4\times10^{5}$ 次联合抽样，最优减投幅度的 95\% 区间为 $[\valHolidayOptimalCutCiLowPct,\ \valHolidayOptimalCutCiHighPct]\%$。令最优减投恰等于实际减投 \valHolidayActualCutPct\% 的临界弹性为 $\theta^{*}=\valHolidayThetaStar$，它落在 $\theta$ 的 95\% 区间 $[\valThetaCiLow,\valThetaCiHigh]$ 之内；若改用问题三、四所用的合并点击--消费弹性 \valPooledGamma，最优减投为 \valHolidayCutAtPooledGammaPct\%，结论由“减投过度”翻转为“减投不足”。因此\textbf{数据不能判定公司的假日减投是过度还是不足}，只能说实际值落在可解释的范围内。同一把尺子用于周几：周一的需求侧点估计给出最优减投 \valMondayOptimalCutPct\%（95\% 区间 $[\valMondayOptimalCutCiLowPct,\ \valMondayOptimalCutCiHighPct]\%$）而实际减投 \valMondayActualCutPct\%，周日为 \valSundayOptimalCutPct\% 对 \valSundayActualCutPct\%——两者与假日一样是“实际减投深于点估计的最优值”，差距分别为 \num{25.7} 与 \num{21.2} 个百分点，而假日为 \num{12.7} 个百分点。周一、周日的临界弹性 $\theta^{*}=\valMondayThetaStar$ 与 \valSundayThetaStar{} 均\emph{落在} $\theta$ 的 95\% 区间之外，因此在这两天“减投过度”的方向比假日更稳健，是投放时间安排中相对最值得复核之处；但“不显著”不等于“效应为零”，本文不把周一/周日的需求侧效应当作 $0$ 处理。

最后必须声明识别限制：式~\eqref{eq:calreg-ctrl} 中的 $\log S_t$ 由公司内生选择（正是本文指出的假日主动减投），$\theta$ 是条件相关弹性而非投放的因果回报弹性；式~\eqref{eq:holiday-ratio} 只在 $\theta$ 可作边际回报解释时成立，其结论应读作参考值而非规范结论（见第~\ref{sec:limits} 节）。

\subsection{注册归因与单元级结构参数}
\label{sec:q1attr}
附件只有日注册总量，为得到各单元每次点击的价值，建立带日历项的归因回归（假设~\ref{as:attr}）
\begin{equation}\label{eq:attr}
R_t=c_0+\sum_{u\in\mathcal U} r_u\,\tilde C_{u,t}+\bm\beta^{\top}\bm D_t+e_t,\qquad
\tilde C_{u,t}=C_{u,t}+\lambda_{\mathrm{ad}}\,\tilde C_{u,t-1},\quad r_u\ge 0,
\end{equation}
其中 $\bm D_t$ 为式~\eqref{eq:calreg} 的日历哑变量，$\tilde C$ 为 adstock 点击（衰减 $\lambda_{\mathrm{ad}}$ 在 $\{0,0.2,0.4,0.6\}$ 中由滚动起点交叉验证的 MAE 选取），系数 $r_u$ 用非负最小二乘\cite{lawson1995}估计。点击数不足 10\,000 或估计值落在合并估计 $\bar r$ 的 $[0.25,4]$ 倍之外的单元使用合并估计（把所有单元点击加总后同样回归得到）。式~\eqref{eq:attr} 选出 $\lambda_{\mathrm{ad}}=\valAttributionLambda$（无显著滞后效应），全样本 $R^2=\valAttributionRsq$，以前 70\% 天数训练、后 30\% 留出的 $R^2=\valAttributionHoldoutRsq$、MAPE \valAttributionHoldoutMape\%（图~\ref{fig:attr}）。截距 $c_0=\valBaselineRegsPerDay$ 人/天是\emph{参照格}（1 月、周三、非假日）的截距，不是全年的非 SEM 基线。全年 \valTotalRegs{} 人的注册按式~\eqref{eq:attr} 恰好分解为三块：截距 $365c_0=\valDecompInterceptRegs$ 人（\valDecompInterceptPct\%）、日历项合计 \valDecompCalendarRegs{} 人（\valDecompCalendarPct\%）与 SEM 点击项 \valDecompSemRegs{} 人（\valDecompSemPct\%），三者相加等于 \valTotalRegs；隐含的非 SEM 日均基线为 \valDecompNonSemPerDay{} 人/天。合并注册率为 \valPooledRegRatePerHundredClicks{} 人/百次点击；\valUnitsWithOwnRate{} 个单元保留自身估计，按下游实际采用的 $r_u$ 归因到 SEM 点击的注册合计 \valAttributedRegs{} 人（占全年注册 \valAttributedRegsSharePct\%，与上面用原始 NNLS 系数的 \valDecompSemRegs{} 人口径不同），平均归因注册成本 \valCostPerAttributedReg{} 元（表~\ref{tab:attr}）。

$r_u$ 是问题二效益分量与问题三、四目标系数的核心结构参数，其不确定性必须与点估计一起报告。表~\ref{tab:attr} 给出每个单元的原始 NNLS 估计、是否触非负边界以及 \valAttributionBoot{} 次移动分块（块长 7 天）残差自助的 95\% 区间：12 个单元中 \valAttributionZeroBoundUnits{} 个的 $\hat r_u$ 被非负约束顶在 0，区间因而单侧退化。识别之所以困难，是因为各单元的点击序列在扣除日历设计后仍高度共线：残差化后点击设计阵的条件数为 \valAttributionConditionNumber。设定稳健性同样显著——去掉月份哑变量后单元注册率最大摆动为 \valAttributionMaxSwingRatio{} 倍。这一不确定性在第~\ref{sec:limits} 节被列为局限，并解释了本文为何以\emph{相对提升}而非绝对注册量作为问题三、四的优越性证据。

单元级点击--消费弹性由活跃日的对数回归估计（假设~\ref{as:resp}）：
\begin{equation}\label{eq:elast}
\log C_{u,t}=a_u+\gamma_u\log S_{u,t}+\sum_{w}\beta_{u,w}D^{w}_t+e_{u,t},
\end{equation}
合并模型（单元固定效应）给出 $\gamma=\valPooledGamma$（SE \valPooledGammaSe）。这里区分两个概念：\emph{回归估计} $\hat\gamma_u$ 与模型\emph{采用值} $\gamma_u$。12 个单元的回归估计为 $\hat\gamma_u\in[\valMinUnitGammaHat,\valMaxUnitGammaHat]$；采用规则是“活跃天数 $n_u\ge 60$ 且 $\hat\gamma_u\in[0.3,0.95]$ 时用自身估计，否则退回合并估计”，其中活跃天数指消费与点击同时为正的天数（表~\ref{tab:attr} 的 $n_u$ 列），与表~\ref{tab:unit-kpis} 的“投放天数”（有记录的天数）不同。据此 \valUnitsWithOwnGamma{} 个单元采用自身估计（$\valMinUnitGamma$--$\valMaxUnitGamma$），另两个退回合并估计 \valPooledGamma：品牌型单元 \valGammaAboveOneUnitId{} 的 $\hat\gamma=\valGammaAboveOneHat$（SE \valGammaAboveOneSe，$n_u=\valGammaAboveOneActiveDays$）\emph{显著大于} 1（$H_0:\gamma=1$ 的 $t=\valGammaAboveOneT$），这与假设~\ref{as:resp} 的凹性相悖，来源是品牌流量的爆发式波动（少数日的点击量跳升数倍），本文按规则将其排除并在第~\ref{sec:limits} 节说明；单元 \valGammaShortUnitId{} 虽有 \valGammaShortDays{} 条记录，但活跃天数只有 \valGammaShortActiveDays{} 天（$\hat\gamma=\valGammaShortHat$），不足以识别。图~\ref{fig:resp} 表明双对数尺度上的线性关系在各单元都成立，$R^2$ 在 \valMinUnitGammaRsq--\valMaxUnitGammaRsq{} 之间（表~\ref{tab:attr}）。

\papertable{tab_attribution_response}{各单元的注册率（式~\eqref{eq:attr}）、归因注册与注册成本，以及弹性估计（式~\eqref{eq:elast}）与采用值}{tab:attr}

\paperfigure[0.8\linewidth]{fig_attribution.pdf}{日新注册数与归因回归~\eqref{eq:attr} 的拟合值（竖虚线右侧为留出期；阴影为法定节假日）。}{fig:attr}

\paperfigure[0.6\linewidth]{fig_response_curves.pdf}{各推广单元日点击量--日消费额散点（双对数）与幂函数响应曲线 $c_u(S/s_u)^{\gamma_u}$；标题中带 * 的单元使用合并弹性。}{fig:resp}

\subsection{合理性小结}
综合以上事实：(i) \valOtherRegionImpSharePct\% 的展现落在上方位之外、只换来 \valOtherRegionClickSharePct\% 的点击，展现结构严重偏向低 CTR 区域（位置效应与创意质量在本数据上不可分离）；(ii) 关键词库长尾严重、四成的词零消费、\valSharedKeywordIds{} 个词跨单元重复投放，缺乏清理与去重；(iii) 预算过度集中于采用弹性最低、CPC 最高的单元，跨单元再分配的反事实收益为 \valQthreeCrossUnitGainPctFeb\%--\valQthreeCrossUnitGainPctAug\%；(iv) 假日、周一与周日的实际减投都深于边际均衡的点估计（分别超出约 \num{12.7}、\num{25.7}、\num{21.2} 个百分点），其中周一、周日的方向判断更稳健，但三者都未被数据显著判定。第 (ii) 点由问题二（分类清理）处理；第 (iv) 点的时间安排由问题三、四的跨日分配处理；第 (iii) 点的跨单元再分配\emph{不}在问题三、四的主模型中执行（假设~\ref{as:budget} 锁定单元预算），只作为反事实与灵敏度情形报告。
